\documentclass[aspectratio=169,11pt]{beamer}
\usepackage[spanish]{babel}
\usepackage[T1]{fontenc}
\usepackage[utf8]{inputenc}
\usepackage{lmodern}
\usepackage{ragged2e}
\usepackage{xcolor}
\usepackage{soulutf8}

\usetheme{Madrid}
\setbeamertemplate{navigation symbols}{}
\setbeamertemplate{footline}[frame number]

\definecolor{azul}{RGB}{33,87,146}
\definecolor{rosa}{RGB}{196,83,138}
\definecolor{verde}{RGB}{74,143,74}
\definecolor{naranja}{RGB}{191,112,33}

\setbeamercolor{structure}{fg=azul}
\setbeamercolor{frametitle}{fg=white,bg=azul}
\setbeamercolor{title}{fg=azul}

\title{Semana 02: Sistematizacion del Tema}
\author{Grupo 6}
\date{CIB12 - Aprendizaje de Maquina y Mineria de Datos}

\begin{document}

\begin{frame}
  \titlepage
\end{frame}

\begin{frame}{Contexto del Tema}
\justifying
\textbf{Linea de investigacion:} Inteligencia Artificial y Comunicacion Inalambrica.

\vspace{0.6em}

\textbf{Tema general:} Museos y espacios culturales.

\vspace{0.6em}

\textbf{Idea central del trabajo:} desarrollar una propuesta de localizacion indoor con aprendizaje profundo y mecanismos de atencion para fortalecer servicios de guia virtual en espacios culturales.
\end{frame}

\begin{frame}{Temas Particulares}
\justifying
\begin{enumerate}
  \item Servicios de guia virtual en museos y espacios culturales.
  \item Servicios de guia virtual con aplicaciones moviles e IoT con senales Wi-Fi/BLE en museos y espacios culturales.
  \item Servicios de guia virtual con deteccion de proximidad y localizacion indoor en museos y espacios culturales.
\end{enumerate}
\end{frame}

\begin{frame}{Temas Especificos}
\justifying
\begin{enumerate}
  \item Baja calidad del servicio de guia virtual en servicios de guia virtual con aplicaciones moviles e IoT con senales Wi-Fi/BLE en museos y espacios culturales.
  \item Baja precision de localizacion indoor en servicios de guia virtual con aplicaciones moviles e IoT con senales Wi-Fi/BLE en museos y espacios culturales.
  \item Alta sensibilidad a la heterogeneidad de dispositivos, inestabilidad del RSSI, efecto multipath y variaciones temporales en servicios de guia virtual con aplicaciones moviles e IoT con senales Wi-Fi/BLE en museos y espacios culturales.
\end{enumerate}
\end{frame}

\begin{frame}{Problema de Estudio}
\justifying
\textbf{Problema de estudio:}

\vspace{0.5em}

\colorbox{yellow!25}{\parbox{0.94\linewidth}{Baja calidad del servicio de guia virtual en servicios de guia virtual con aplicaciones moviles e IoT con senales Wi-Fi/BLE en museos y espacios culturales, 2026.}}

\vspace{0.9em}

\textbf{Interpretacion:}
\begin{itemize}
  \item El servicio depende de una localizacion indoor suficientemente precisa.
  \item El entorno presenta ruido, multipath y heterogeneidad de dispositivos.
  \item Esto afecta la activacion contextual correcta de la guia virtual.
\end{itemize}
\end{frame}

\begin{frame}{Titulos y Variables}
\justifying
\textbf{Titulo preliminar:}

\vspace{0.4em}

Calidad del servicio de guia virtual en servicios de guia virtual con aplicaciones moviles e IoT con senales Wi-Fi/BLE en museos y espacios culturales, 2026.

\vspace{0.7em}

\textbf{Titulo tentativo:}

\vspace{0.4em}

Modelo de aprendizaje profundo con mecanismos de atencion para mejorar la calidad del servicio de guia virtual en servicios de guia virtual con aplicaciones moviles e IoT con senales Wi-Fi/BLE en museos y espacios culturales, 2026.

\vspace{0.7em}

\textbf{Variable independiente:} Modelo de aprendizaje profundo con mecanismos de atencion.

\textbf{Variable dependiente:} Calidad del servicio de guia virtual.

\textbf{Objeto de estudio:} Servicios de guia virtual.
\end{frame}

\begin{frame}{Alcances y Verificacion}
\justifying
\textbf{Alcance espacial:} Museos y espacios culturales.

\textbf{Alcance temporal:} 2026.

\vspace{0.8em}

\textbf{Verificacion frente a las diapositivas:}
\begin{itemize}
  \item Se definio linea de investigacion.
  \item Se formulo tema general.
  \item Se propusieron tres temas particulares.
  \item Se construyeron temas especificos.
  \item Se definio problema de estudio.
  \item Se obtuvo titulo preliminar y titulo tentativo.
  \item Se delimitaron variables, objeto de estudio y alcances.
\end{itemize}
\end{frame}

\begin{frame}{Siguiente Entregable}
\justifying
\textbf{Entrega actual:} Sistematizacion del tema de investigacion.

\vspace{0.7em}

\textbf{Siguiente entregable:} Antecedentes bibliograficos.

\vspace{0.7em}

\textbf{Trabajo inmediato del grupo:}
\begin{itemize}
  \item seleccionar antecedentes directamente relacionados con indoor localization;
  \item resumir autores, año, metodo y hallazgo;
  \item identificar vacios de la literatura que justifiquen el tema.
\end{itemize}
\end{frame}

\end{document}
